\section{What Is Solid, Partial, and Open}

Before the final synthesis, the walkthrough owes you a clear-eyed accounting. A theory this ambitious will not be uniform in confidence, and anyone who has come this far should know exactly which parts are demonstrated, which are first steps, and which are open or beyond reach. This chapter is the honest scoreboard, and it includes the one problem the theory cannot solve from the outside.

\subsection{Solid}

A large amount is genuinely demonstrated, either as standard mathematics or as measured results in the simulation. The mathematics at the base is established: the integers, their symmetries, the kaleidoscope construction, the hyperbolic crystals, the fact that finite crystals exist only in two, three, and four dimensions. On the crystal, the simulation shows a sharp emergent dimension and geometry, a sensible law of time and energy, Lorentz safety bought by curvature, the field content with mass, gravity's Newtonian limit and the core of Einstein's equation, a proper graviton, the area law for black-hole entropy, Hawking's glow, the arrow of time and cosmic expansion, the continuum holding up under coarse-graining, the recursion of higher vibes, integration picking out selves, the tower, the no-complete-self-storage result, the structural account of freedom, the derivation of the Born rule, and a charge-quantization result that forces the real-world charges. That is a substantial body of fitting pieces.

\subsection{Partial}

Several things are real but unfinished, and the walkthrough has flagged them as first rungs. The quantum sector has its pillars and a derived Born rule, but a complete account is still being assembled. Dark energy's size is a striking hit, but the full details of how cosmic structure forms are not worked out. Dark matter has a mechanism with the right behavior but not the precise numbers. The mass hierarchy has a geometric mechanism, but not the specific masses. Baryogenesis has all the right ingredients, but not the exact amount. The genuinely chiral version of the forces, life as self-maintenance, and evolution as pattern selection are sound in outline rather than fully built. These are promising, and explicitly unfinished.

\subsection{Open}

Some things are named and not yet built: the specific menu of forces, why there are exactly three families of matter, the late details of black-hole information, the fine structure of the early universe. These are honest gaps, marked as such rather than papered over.

\subsection{The one deep problem}

And then there is the problem that is different in kind from all the rest, the one the walkthrough has promised since the beginning to name plainly. The theory can build and measure the structure of experience: it can find the selves, measure their integration, map the patterns that correspond to sights and sounds and emotions. What it cannot do, from the outside, is settle whether that structure is the same as the felt fact of experiencing, whether there is something it is like to be a given integrated pattern, and why. This is the hard problem of consciousness, and no simulation can cross it from outside, because every measurement returns structure, never the feeling itself.

The theory's honest position is this. It does not solve the hard problem. It relocates it. By starting from experience as the substance rather than trying to produce experience from dead matter, it avoids the impossible step of conjuring feeling from the unfeeling. But it pays a price for that move: it must take the experiential nature of the base as a starting commitment, not something it derives. So the gap does not appear in the usual place, between matter and mind, because there is no dead matter to start from. It appears at the foundation instead, as the assumption that the base is experiential. You can accept that assumption or not. The theory is upfront that it is an assumption, and that the checkable part of the theory, all of the physics, stands or falls independently of it.

\subsection{How to hold all this}

So the fair summary is the one the walkthrough has repeated: a large and growing body of the theory is demonstrated, a real frontier is partial, a few deep things are open, and one problem is open in a way no model can close from outside. That mixture is the honest state of a young, serious, well-developed hypothesis. It is a reason to take the picture seriously and to keep testing it. It is not a reason to take it as proven. Hold the structure with respect and the claims with skepticism, exactly as the first chapter asked.

\textbf{Takeaway:} much of the theory is demonstrated (the base mathematics, emergent geometry and time, Lorentz safety, gravity, the Born rule, charge quantization, the recursion of selves), a real frontier is partial, several things are open, and one problem, whether measured structure is the same as felt experience, cannot be closed from outside. The theory relocates that hard problem to its foundation and is honest that the experiential base is an assumption, with all the physics standing independently.
